\section{Supportable hardware}
\label{sec-supportable-hardware}
Any interface to a robot that has a function like
\fbox{\texttt{q = interface.getQ()}} to get the current joint positions
and a function to send velocity commands like
\fbox{\texttt{interface.sendQd(qd)}}
can be made to comply with SMC.
Along with this, one needs to have a URDF of the robot
or manually construct the Pinocchio model.

The requirements are codified contract via 
\fbox{\texttt{AbstractRealRobotManager}}'s 
abstract methods.
This class inherits from AbstractRobotManager,
and exist simply because it has additional methods not needed in simulation
(for simulation we use AbstractSimulatedRobotManager which integrates
joint positions instead of reading them from the real robot).


Here are methods which need to implemented
in the \fbox{\texttt{RealXYZRobotManager}} to get SMC to control a real robot XYZ.
\begin{itemize}
		\item \fbox{\texttt{connectToRobot()}}
		\item \fbox{\texttt{\_updateQ()}}
		\item \fbox{\texttt{sendVelocityCommandToReal()}}
		\item \fbox{\texttt{setInitialPose()}}
		\item \fbox{\texttt{stopRobot()}}
\end{itemize}

The following functions are also made abstract,
but could trivially implemented (ex. just an empty return)
if not available or required:
\begin{itemize}
		\item \fbox{\texttt{connectToGripper()}}
		\item \fbox{\texttt{\_updateV()}}
		\item \fbox{\texttt{setFreedrive()}}
		\item \fbox{\texttt{unSetFreedrive()}}
\end{itemize}
